\subsubsection{tikz}

\def\angle{70} % launch angle
\def\v{6}   % initial velocity
\def\g{9.8} % gravitational acceleration
\def\r{0.1} % ball radius

\begin{center}
\resizebox{0.5\textwidth}{!}{
\begin{tikzpicture}[
    scale=2.5,
    axis/.style={-stealth, thick},
    every node/.style={font=\small},
]


    \pgfmathsetmacro{\tGround}{{2*\v*sin(\angle)/\g}}
    \newcommand{\Ballx}[1]{{\v*#1*\tGround*cos(\angle)}}
    \newcommand{\Bally}[1]{{-0.5*\g*#1*#1*\tGround*\tGround+\v*#1*\tGround*sin(\angle)}}
    
    % Axis
    \draw[axis] (-0.3,0) -- ($1.2*(\Ballx{1},0)$) node[right] {$x$};
    \draw[axis] (0,-0.3) -- ($1.2*(0,\Bally{0.5})$) node[above] {$y$};

    % Angle
    \filldraw[opacity=0.5] (\r*1.6,0) arc(0:\angle:\r*1.6)node[opacity=1,midway,above right]{$\theta$} -- (0,0) -- cycle;

    % Ball
    \newcommand{\DrawBall}[2][black]{
        \filldraw[ball color=#1!50!] (\Ballx{#2},\Bally{#2}) circle(\r);
        \draw[-stealth,#1] (\Ballx{#2},\Bally{#2})
            --++({0.1*\v*cos(\angle)}, {-0.1*\g*#2*\tGround + 0.1*\v*sin(\angle)}); % velocity
    }

    \foreach \t in {0,0.2,0.8,1} {
        \DrawBall[black]{\t}
    }
    
    \draw[densely dotted] (\Ballx{0.2},0) node[below] {$x=v_0t\cos\theta$}
        -- (\Ballx{0.2},\Bally{0.2})
        -- (0,\Bally{0.2}) node[left] {$\displaystyle y=-\frac{1}{2}gt^2+v_0t\sin\theta$};
    
    \DrawBall[red]{0.5}
    
    % Trajectory
    \draw[samples=256,densely dotted,domain=0:1,variable=\t] plot(\Ballx{\t},\Bally{\t});
    % 关键：在图片正下方添加说明文字（相对图片居中）
    \node[align=center,yshift=-3.5em,font=\small] at ($(\Ballx{1},0)!0.5!(0,0)$)
        {抛体运动轨迹示意图};

\end{tikzpicture}
 }

\end{center}



\subsection*{cal}

\begin{center}
\resizebox{0.5\textwidth}{!}{
\begin{tikzpicture}
  \calendar
  [
    dates=\year-\month-\day+-25 to \year-\month-\day+25,
    week list,inner sep=2pt,month label above centered,
    month text=\textit{\%mt \%y0}
  ]
  if (at least=\year-\month-\day) {}
    else [nodes={strike out,draw}]
  if (at most=\year-\month-\day+7)
    [green!50!black]
  if (between=\year-\month-\day+8 and \year-\month-\day+10)
    [red]
  if (Sunday)
    [gray,nodes={draw=none}]
  ;
\end{tikzpicture}
}
\end{center}


\subsubsection{polar axes}
\begin{center}
\resizebox{0.8\textwidth}{!}{
\tikz[baseline] \datavisualization [
  scientific axes={clean},
  x axis={attribute=angle, ticks={minor steps between steps=4}},
  y axis={attribute=radius, ticks={some, style=red!80!black}},
  all axes=grid,
  visualize as smooth line=sin]
  data [format=function] {
    var t : interval [-3:3];
    func angle = exp(\value t);
    func radius = \value{t}*\value{t};
  };
\qquad
\tikz[baseline] \datavisualization [
  scientific polar axes={right half clockwise, clean},
  angle axis={logarithmic,
    ticks={
      minor steps between steps=8,
      major also at/.list={2,3,4,5,15,20}}},
  radius axis={ticks={some, style=red!80!black}},
  all axes=grid,
  visualize as smooth line=sin]
  data [format=function] {
    var t : interval [-3:3];
    func angle = exp(\value t);
    func radius = \value{t}*\value{t};
  };
}
\end{center}

\subsubsection*{tikzdatavisualizationset}
\tikzdatavisualizationset {
  example visualization/.style={
    scientific axes=clean,
    y axis={ticks={style={
          /pgf/number format/fixed,
          /pgf/number format/fixed zerofill,
          /pgf/number format/precision=2}}},
    x axis={ticks={tick suffix=${}^\circ$}},
    1={label in legend={text=$\frac{1}{6}\sin 11x$}},
    2={label in legend={text=$\frac{1}{7}\sin 12x$}},
    3={label in legend={text=$\frac{1}{8}\sin 13x$}},
    4={label in legend={text=$\frac{1}{9}\sin 14x$}},
    5={label in legend={text=$\frac{1}{10}\sin 15x$}},
    6={label in legend={text=$\frac{1}{11}\sin 16x$}},
    7={label in legend={text=$\frac{1}{12}\sin 17x$}},
    8={label in legend={text=$\frac{1}{13}\sin 18x$}}
  }
}

% 定义共用数据集 sin functions
\tikz \datavisualization data group {sin functions} = {
  data [format=function] {
    var set : {1,...,8};
    var x : interval [0:50];
    func y = sin(\value x * (\value{set}+10))/(\value{set}+5);
  }
};

\subsubsection*{vary thickness}
\begin{center}
\resizebox{0.8\textwidth}{!}{
\begin{tikzpicture}
\datavisualization [
  visualize as smooth line/.list={1,2,3,4,5,6,7,8},
  example visualization,
  style sheet=vary thickness
]
data group {sin functions};
\end{tikzpicture}
}
\end{center}
